\scp{Terms in Wallacean subgroups glossed ‘bandicoot’}{14-05}
Bandicoots (\it{Peramelidae spp.}) are not found in most
Wallacean subgroups,\footnote{
		\citet[197]{Ellen1993}, in his detailed study of animals in Maluku observes,
		“The bandicoot or ‘marsupial rat’ is rare in the Moluccas.” \citeauthor{Ellen1993}’s
		comment is framed as an objection to \citeauthor{Blust1982}'s (\citeyear[241]{Blust1982}) claims
		about widespread glosses of ‘bandicoot’ throughout the region.}
and in some subgroups their range of habitat does not encompass the whole subgroup.
Terms in Wallacean subgroups glossed ‘bandicoot’ (that are not
actually terms for ‘cuscus’ erroneously glossed as ‘bandicoot’
discussed in \srf{14-04-01}), are shown in \trf{14-14} (see appendix \ref{ch:SurMar} for sources),
along with known terms from SHWNG languages.
Note that the terms given in \trf{14-14} refer to several different species
of bandicoot (see \citealp[268]{Schapper2011}).
%(STB = \tsc{Seram-Tanimbar-Bomberai}).

As can be seen from the forms glossed ‘bandicoot’ in \trf{14-14},
there is great diversity of forms,
suggesting that terms for the elusive bandicoot developed independently
from region to region.
There is consistency across some Aru languages,
such that \citet[267]{Schapper2011} reconstructs **kawaran ‘bandicoot’
for Manombai, Lorang, and West Kola.
Southwest Tarangan \it{man} is clearly anomalous within Aru for
words glossed ‘bandicoot?’, and is also anomalous for
reflexes of Proto-Aru **medaR ‘cuscus’ (see \srf{14-04-03}).
Anomalous Southwest Tarangan \it{man} ‘bandicoot’ is the only
possible reflex of **mans(aə)r outside of Oceanic languages that
lends support to \citeauthor{Blust1982}’s (\citeyear{Blust1982,Blust2012b})
insistence of the gloss of ‘bandicoot’ for this etymon
for PCEMP.

\begin{table}[p]\centering\tcp{Wallacean words glossed `bandicoot' or similar}{14-14}
%\begin{adjustbox}{width=\textwidth}
	\begin{threeparttable}\st{3.5pt}
	\begin{tabular}{l|lll}\lsptoprule
Language	&	Form	&	local gl.	&	Node	\\	\midrule
S. Nuaulu	&	imanona	&	long-nosed marsupial\su{†}	&	\tsc{AS}	\\
Manusela\su{‡}	&	mape-a	&	Seram long-nosed bandicoot,	&\tsc{AS}\\
		&		&	\it{Rhynchomeles prattorum}	&		\\
	&	mabaya	&	bandicoot	&	\tsc{AS}	\\	\midrule
Kei Besar	&	helé	&	bandicoot, \it{Echymipera rufescens}	&	STB	\\
Kei Kecil	&	halee	&	mouse-like marsupial	&	STB	\\	\midrule
Ujir	&	koa	&	bandicoot, \it{Echymipera rufescens}	&	\tsc{Aru}	\\
East Kola	&	koyi	&	bandicoot, \it{Echymipera rufescens}	&	\tsc{Aru}	\\
West Kola	&	kawaran	&	bandicoot, \it{Echymipera rufescens}	&	\tsc{Aru}	\\
Batuley	&	koyi	&	bandicoot, \it{Echymipera rufescens}	&	\tsc{Aru}	\\
				&	karsirem	&	\it{tikus tanah} (`earth rat') which&	\tsc{Aru}	\\
				&						&	lives in cliff holes	&	\tsc{Aru}	\\
Dobel	&	ʔosi	&	bandicoot, \it{Echymipera rufescens}	&	\tsc{Aru}	\\
Koba	&	ŋarukwabala	&	bandicoot, \it{Echymipera rufescens}	&	\tsc{Aru}	\\
Lorang	&	kagwaran	&	bandicoot, \it{Echymipera rufescens}	&	\tsc{Aru}	\\
Manombai	&	kagaran	&	bandicoot, \it{Echymipera rufescens}	&	\tsc{Aru}	\\
SW Tarangan	&	man	&	bandicoot?	&	\tsc{Aru}	\\
	&	karsiram	&	a kind of marsupial	&	\tsc{Aru}	\\	\midrule
Matbat	&	isɔ¹²k	&	bandicoot, \it{Echymipera kalubu}	&	SHWNG	\\
Biak	&	karau (sop)	&	bandicoot, \it{Echymipera kalubu}	&	SHWNG	\\
Ma{\Q}ya	&	kalubu	&	bandicoot, \it{Echymipera kalubu}	&	SHWNG	\\
Ansus	&	pimunau	&	bandicoot (unspecified)	&	SHWNG	\\
Moor	&	atíata	&	bandicoot (unspecified)	&	SHWNG	\\
Ambai	&	kafa	&	bandicoot (generic)	&	SHWNG	\\
\lspbottomrule
	\end{tabular}
		\begin{tablenotes}\tnsize
			\item[†]	{\citet{MatokeBolton2005} give \it{imanona} ‘tikus besar’ = ‘large rat/mouse’.
								The more specific gloss is from \citet{Bolton-Grimes1990}.
								Roy Ellen (pers. comm. July 2024) doubts that \it{imanona} is
								`bandicoot'. He states ``I collected lots of non-marsupial rats that were
								confidently identified by people as \it{imanona}.
								However, descriptions of \it{tui-tui} come close to what we know of \it{R.
								prattorum}, though of course \it{tui-tui} may be a synonym of other
								terms that I never discovered.''}
			\item[‡]	{\it{mape-a} with the more specific gloss is from \citet{MacDonald2015}. Alistair A.
								MacDonald is a zoologist specializing in mammals of Indonesia
								and we therefore have extremely high confidence in this gloss.
								\it{Mabaya}, with a less specific gloss in \citet[268]{Schapper2011}
								is from \citet{Flannery1995}.}
		\end{tablenotes}
	\end{threeparttable}
%\end{adjustbox}
\end{table}